% --- MODULE 1: Algorithm Analysis & Amortization ---
% --- LECTURES 1, 2, 3, 9 & CLRS 16 ---

\section{Module 1: Algorithm Analysis \& Amortization}

\subsection{Core Concepts}
\begin{description}
    \item[Worst-Case Analysis] (Lec 2) Measures the maximum number of operations an algorithm will perform for any input of size $n$. This is the standard for analysis in this course.
    \item[Amortized Analysis] (Lec 9, CLRS 16) Analyzes a \textit{sequence} of operations to find the average cost per operation in the worst case. It is \textbf{not} the same as average-case analysis.
\end{description}

\subsection{Asymptotic Notations (Lec 3)}
\noindent For a given time function $T(n) \ge 0$:
\begin{description}
    \item[$O(f(n))$ (Big-Oh)] (Upper Bound) $\exists$ positive constants $c, n_0$ such that $T(n) \le c \cdot f(n)$ for all $n \ge n_0$.
    \item[$\Omega(f(n))$ (Big-Omega)] (Lower Bound) $\exists$ positive constants $c, n_0$ such that $T(n) \ge c \cdot f(n)$ for all $n \ge n_0$.
    \item[$\Theta(f(n))$ (Big-Theta)] (Tight Bound) $\exists$ positive constants $c_1, c_2, n_0$ s.t. $c_1 f(n) \le T(n) \le c_2 f(n)$ for all $n \ge n_0$. (This means $T(n)$ is $O(f(n))$ and $\Omega(f(n))$).
\end{description}

\subsection{Common Growth Rates (Lec 3)}
$O(1) \subset O(\log n) \subset O(n) \subset O(n \log n) \subset O(n^2) \subset O(n^k) \subset O(2^n) \subset O(n!)$

\subsection{Amortized Analysis Methods (Lec 9, CLRS 16)}
\subsubsection{1. Aggregate Analysis (Lec 9)}
Show that for a sequence of $n$ operations, the total cost $T(n)$ is $O(f(n))$. The amortized cost per operation is then $T(n)/n$.
\begin{itemize}
    \item \textbf{Binary Counter ($n$ increments):} Total cost = $\sum_{i=0}^{\log n} \frac{n}{2^i} \approx 2n = O(n)$. Amortized cost = $\mathbf{O(1)}$.
    \item \textbf{Dynamic Array (Push):} Total cost for $n$ pushes = $O(n)$ (for $n$ pushes) + $O(n)$ (for all copying). Amortized cost = $\mathbf{O(1)}$.
\end{itemize}

\subsubsection{2. Accounting Method (CLRS 16.2)}
Assign different \textit{amortized costs} ($\hat{c}$) to operations.
\begin{itemize}
    \item If $\hat{c} > c$ (actual cost), the extra charge is stored as \textbf{credit} on the data structure.
    \item If $\hat{c} < c$, the deficit is paid by stored credit.
    \item \textbf{Rule:} The total credit must \textbf{never be negative}.
    \item \textbf{Dynamic Array (Push):}
    \begin{itemize}
        \item Let \textbf{Amortized Cost $\hat{c} = 3$}.
        \item \textbf{Case 1: No resize (actual cost $c=1$)}. We pay 3. 1 pays for the push, 2 are stored as credit.
        \item \textbf{Case 2: Resize (actual cost $c=n+1$)}. Occurs when $n-1$ items are present. There must be $(n-1)$ items, each with 2 units of credit, for a total of $2(n-1)$ credit. We use $n$ of this credit to pay for the copy, and 1 to pay for the push.
        \item We use $n+1$ credit, but have $2n-2$. Credit remains non-negative. Amortized cost is $\mathbf{O(1)}$.
    \end{itemize}
\end{itemize}

\subsubsection{3. Potential Method (CLRS 16.3)}
Define a potential function $\Phi$ that maps the data structure $D$ to a number.
\begin{itemize}
    \item $\Phi(D_0) = 0$ and $\Phi(D_i) \ge 0$ for all $i$.
    \item Amortized cost: $\hat{c_i} = c_i + \Phi(D_i) - \Phi(D_{i-1})$ (actual cost + change in potential).
    \item \textbf{Dynamic Array (Push):}
    \begin{itemize}
        \item Define $\Phi(T) = 2 \cdot \text{num}(T) - \text{size}(T)$. (Assumes $\text{size}/2 < \text{num} \le \text{size}$).
        \item \textbf{Case 1: No resize ($c_i=1$)}. $\text{num}_i = \text{num}_{i-1}+1$, $\text{size}_i = \text{size}_{i-1}$.
        \item $\Delta\Phi = \Phi(D_i) - \Phi(D_{i-1}) = (2 \cdot \text{num}_i - \text{size}_i) - (2 \cdot (\text{num}_i-1) - \text{size}_i) = 2$.
        \item $\hat{c_i} = c_i + \Delta\Phi = 1 + 2 = 3$.
        \item \textbf{Case 2: Resize ($c_i=n$)}. $\text{num}_i = n$, $\text{size}_i = 2(n-1)$. $\text{num}_{i-1} = n-1$, $\text{size}_{i-1} = n-1$.
        \item $\Delta\Phi = (2n - 2(n-1)) - (2(n-1) - (n-1)) = 2 - (n-1) = 3 - n$.
        \item $\hat{c_i} = c_i + \Delta\Phi = n + (3 - n) = 3$.
        \item In both cases, the amortized cost is $\mathbf{O(1)}$.
    \end{itemize}
\end{itemize}